%-------------------------
% Cover Letter - Waymo Software Engineer, Perception Evaluation
% Author: Ajay Venkatesh
%-------------------------

\documentclass[letterpaper,11pt]{article}

\usepackage[top=0.8in, bottom=0.8in, left=0.9in, right=0.9in]{geometry}
\usepackage[hidelinks]{hyperref}
\usepackage{lmodern}
\usepackage{parskip}
\usepackage{microtype}

\setlength{\parskip}{6pt}
\setlength{\parindent}{0pt}

\begin{document}

\begin{flushleft}
\textbf{\large Ajay Venkatesh} \\
Brooklyn, NY \\
+1 516 585 9013 \\
\href{mailto:av3855@nyu.edu}{av3855@nyu.edu}
\end{flushleft}

\vspace{10pt}

May 13, 2026

\vspace{6pt}

Hiring Manager \\
Waymo

\vspace{10pt}

Dear Hiring Manager,

My name is Ajay Venkatesh --- finishing my M.S. in Computer Engineering at NYU, with several years of production software experience: a few years at PwC in Bengaluru, a summer SDE internship at Amazon, and ongoing on-campus work at NYU's Novel AI Technologies. I'm writing about the Software Engineer, Perception Evaluation role. The work the team describes --- building the tools and pipelines that quantify whether a Perception system is doing what it claims --- is the part of ML engineering I find most interesting, and the part where my background lines up most cleanly.

\textbf{Python} is my daily driver for both backend services and ML pipelines --- I've shipped data preprocessing, model integration, evaluation tooling, and end-to-end ETL workflows across multiple projects. \textbf{C++} is where I cut my teeth in graduate coursework (Computer Architecture, System Design, Data Structures) and where I keep going back when an abstraction stops being honest. Clean, well-tested code is the baseline: at Amazon I shipped backend services in Java with unit/integration testing via \textbf{Mockito}, low-level design patterns, and code-review discipline, and the same posture carries over to my Python work.

On pipelines and evaluation infrastructure specifically: at Amazon I built a high-throughput, event-driven processing system on \textbf{AWS ECS Fargate, SQS, SNS} with auto-scaling clusters and resilient \textbf{Infrastructure-as-Code} on \textbf{AWS CDK} --- the same shape of work as a perception evaluation pipeline: scalable jobs, reliability monitoring, and the discipline of catching regressions before they ship. I also deployed a \textbf{Retrieval-Augmented Generation} pipeline with an \textbf{agentic workflow} that triages operational incidents by reasoning over logs --- the practical engineering side of "analyze results, troubleshoot issues, and surface findings" the role calls out.

On large-scale data and analysis: my Big Data project built a \textbf{PySpark / HDFS} pipeline over millions of flight records, training Random Forest and Gradient Boosted Tree models with rigorous EDA. I'm fluent in \textbf{Pandas, NumPy,} and \textbf{Jupyter / Colab} for data analysis, and I've used \textbf{Power BI}, \textbf{AWS QuickSight}, and Matplotlib to turn raw metrics into actionable insights for stakeholders --- the "analyze and report findings" part of the role isn't a separate job for me, it's where the work usually lands.

What pulls me toward Perception Evaluation specifically is the stakes. Perception is the layer where the Waymo Driver's mistakes start being expensive --- in real units, on real roads, with people inside the vehicle. The team's job is to be the conscience of that stack: to surface what the system actually does, not what its top-line metric claims it does. That's the engineering culture I want to be inside of.

I'd welcome the chance to discuss further. Thanks for considering my application.

\vspace{12pt}

Sincerely, \\
Ajay Venkatesh

\end{document}
